\newpage % Rozdziały zaczynamy od nowej strony.
\section{Related work}

In this section, we review the concepts and methods underlying both tasks
addressed in this thesis: knowledge graphs as structured medical knowledge,
graph neural networks for representation learning on heterogeneous graphs,
and Hierarchical Correlation Reconstruction (HCR) as a statistical complement
to learned graph representations.
We begin with graph-based knowledge representation, then describe deep
learning on graphs, and finally introduce HCR.

\subsection{Knowledge graphs}
\label{subsec:related-kg}

Knowledge graphs represent structured knowledge through entities and the
relations between them.
Formally, a knowledge graph $G$ is a set of triples

\begin{definition}[Knowledge Graph]
\label{def:kg}
$$G \subseteq \mathcal{V} \times \mathcal{R} \times \mathcal{V},$$
where $\mathcal{V}$ denotes the set of entities (nodes) and $\mathcal{R}$
the set of relation types (edge labels). Each triple $(h, r, t) \in G$
corresponds to a directed, labeled edge $h \xrightarrow{r} t$, where $h$
and $t$ are referred to as the head and tail entity, respectively
\textcite{hogan2021knowledge}.
\end{definition}

The key benefit of knowledge graphs is that they provide an interlinked
representation of knowledge in a human- and machine-understandable format
(\textcite{barrasa2023building}).
In biomedical applications, entities may represent drugs, diseases, genes,
proteins, mechanisms, patient characteristics, or clinical outcomes, while
edges describe typed relationships between them.
Such graphs are naturally heterogeneous, because both nodes 
and relations can carry different semantic types.
This motivates the heterogeneous-graph formulation:

\begin{definition}[Heterogeneous Graph]
\label{def:heterokg}
A heterogeneous graph extends Definition~\ref{def:kg} with explicit type
information and is defined as a quadruple
$$G = (V, E, \tau, \phi),$$
where $V$ is the set of nodes, $E \subseteq V \times V$ is the set of
edges, and $\tau: V \rightarrow \mathcal{T}_V$, $\phi: E \rightarrow
\mathcal{T}_E$ are type-mapping functions assigning each node and edge a
type from finite type sets $\mathcal{T}_V$ and $\mathcal{T}_E$
\cite{hu2020heterogeneous}. For an edge $e = (u, v) \in E$, the triple
$\bigl(\tau(u), \phi(e), \tau(v)\bigr)$ defines its schema-level relation,
which corresponds directly to the \texttt{edge\_type} representation used
in the \texttt{HeteroData} object in this thesis.
\end{definition}

Knowledge graphs have been widely applied in biomedical research,
particularly for integrating information from multiple sources and
supporting drug discovery (\textcite{maclean2021knowledge}).
However, many biomedical knowledge graphs primarily encode \emph{associations}
extracted from literature and curated databases, not
patient-specific empirical realizations of the mechanisms they describe.
Popular graph-based methods therefore often ``highlight biological
associations, failing to model causal relationships in complex dynamic
biological systems'' \cite{maclean2021knowledge}.
Recent work has increasingly considered integrating knowledge graphs with
causal or mechanistic representations
\cite{toonsi2025causal,lyu2023causal,mesinovic2026causalgnn}.

Beyond biomedicine, knowledge graphs are also combined with empirical
interaction data in recommendation systems, where models such as KGAT
propagate representations over a collaborative knowledge graph built from
user--item interactions and external knowledge \cite{wang2019kgat}.
The Last-FM benchmark from this line of work is used as a cross-domain
proxy in link prediction task.
The reliability of this benchmark was subsequently questioned by
Shehzad et al.\ \cite{shehzad2025reproducibility} in a reproducibility
study of intent-aware recommendation models; the corrected Last-FM* split
used in this thesis is described in Section~\ref{sec:models-benchmark}.

